\subsection{Integração com o Banco de Dados}
\label{sec:banco-dados}

A persistência do módulo foi implementada com SQLAlchemy sobre PostgreSQL, 
seguindo o padrão já adotado pelo sistema. O mapeamento foi feito em estilo 
declarativo, com uma tabela principal (\texttt{solicitations}) e duas tabelas 
de referência (\texttt{solicitation\_status} e \texttt{solicitation\_category}), 
que funcionam como \textit{lookup tables} para os estados e os tipos de 
solicitação. A escolha de manter tabelas de referência, em vez de armazenar 
os valores como texto livre, garante integridade referencial no banco e 
permite a evolução do domínio sem migrations estruturais a cada novo valor.
O Quadro~\ref{qua:modelo-solicitacao} descreve as tabelas do módulo, suas 
colunas e restrições.

\begin{quadro}[htbp]
\centering
\caption[Modelo de dados do módulo de solicitações]{Modelo de dados do módulo de solicitações.}
\label{qua:modelo-solicitacao}
\footnotesize
\setlength{\tabcolsep}{3pt}
\resizebox{\linewidth}{!}{%
\begin{tabular}{p{2.3cm}p{2.0cm}p{3.1cm}p{3.7cm}}
\toprule
\textbf{Tabela} & \textbf{Coluna} & \textbf{Tipo / Restrição} & \textbf{Observação} \\
\midrule
\texttt{solicitations} & \texttt{id} & Inteiro, PK & Identificador da solicitação \\
 & \texttt{id\_user} & FK $\rightarrow$ \texttt{users}, obrigatório & Solicitante \\
 & \texttt{id\_status} & FK $\rightarrow$ \texttt{solicitation\_status}, padrão \texttt{CREATED} & Estado atual \\
 & \texttt{id\_category} & FK $\rightarrow$ \texttt{solicitation\_category}, opcional & Tipo de solicitação; nulo em rascunho \\
 & \texttt{specialist\_type\_id} & FK $\rightarrow$ \texttt{users\_specialist\_types}, opcional & Categoria de destino \\
 & \texttt{title} & Texto(128), obrigatório & Título \\
 & \texttt{description} & Texto(640), obrigatório & Justificativa ou descrição \\
 & \texttt{file\_path} & Texto(256), opcional & Caminho do documento anexado \\
 & \texttt{created\_at} & DataHora, padrão \texttt{agora()} & Criação \\
 & \texttt{last\_update} & DataHora, atualizado em alteração & Última modificação \\
\midrule
\texttt{solicitation status} & \texttt{id}, \texttt{name} & PK, Texto(20) único & \texttt{CREATED}, \texttt{PENDING}, \texttt{APPROVED}, \texttt{REJECTED} \\
\texttt{solicitation category} & \texttt{id}, \texttt{name} & PK, Texto(20) único & \texttt{CHANGE\_CATEGORY}, \texttt{RESOURCE}, \texttt{UPDATE\_POINTS} \\
\texttt{solicitation reviews} & \texttt{id\_solicitation} & FK $\rightarrow$ \texttt{solicitations}, cascata & Decisões e pareceres das análises \\
\bottomrule
\end{tabular}%
}
\fonte{Elaboração própria (2026).}
\end{quadro}

O Algoritmo~\ref{alg:modelo-solicitation} descreve, em pseudocódigo, a 
estrutura da entidade principal, com seus atributos, restrições, índices 
e relacionamentos, e as enumerações que definem os estados e os tipos de 
solicitação.

\begin{algorithm}[htbp]
\caption{Estrutura da entidade \texttt{Solicitation} e enumerações de domínio (pseudocódigo).}
\label{alg:modelo-solicitation}
\begin{verbatim}
ENUMERAÇÃO SolicitationStatus
    CREATED  = 1    // rascunho, em elaboração
    PENDING  = 2    // submetida, aguardando análise
    APPROVED = 3    // deferida
    REJECTED = 4    // indeferida

ENUMERAÇÃO SolicitationCategory
    CHANGE_CATEGORY = 1    // alteração de categoria
    RESOURCE        = 2    // recurso / complemento cadastral
    UPDATE_POINTS   = 3    // atualização de pontuação

ENTIDADE Solicitation
    TABELA: solicitations

    ATRIBUTOS
        id                 : Inteiro, Chave Primária
        id_user            : Inteiro, FK -> User, obrigatório
        id_status          : Inteiro, FK -> SolicitationStatus,
                             padrão CREATED
        id_category        : Inteiro, FK -> SolicitationCategory,
                             opcional (discriminador de tipo)
        specialist_type_id : Inteiro, FK -> UserSpecialistType,
                             opcional (categoria de destino)
        title              : Texto(128), obrigatório
        description        : Texto(640), obrigatório
        file_path          : Texto(256), opcional
        created_at         : DataHora, padrão agora()
        last_update        : DataHora, atualizado em cada alteração

    RELACIONAMENTOS
        user    : User                 (muitos-para-um)
        reviews : Lista<Review>        (um-para-muitos, cascata)

    ÍNDICES
        idx_solicitation_user     -> (id_user)
        idx_solicitation_status   -> (id_status)
        idx_solicitation_created  -> (created_at)
\end{verbatim}
\end{algorithm}

A decisão de modelagem central é a \textbf{abordagem polimórfica em tabela 
única}: os três tipos de solicitação (\texttt{CHANGE\_CATEGORY}, 
\texttt{RESOURCE} e \texttt{UPDATE\_POINTS}) compartilham a mesma tabela 
\texttt{solicitations}, discriminados pelo campo \texttt{id\_category}. 
Essa escolha evita a fragmentação em três tabelas com colunas comuns 
replicadas (título, descrição, autor, estado, timestamps) e concentra a 
lógica de listagem, filtragem e paginação em uma única consulta, 
independentemente do tipo. Campos específicos de um tipo, como 
\texttt{specialist\_type\_id}, são opcionais e populados apenas quando 
aplicáveis: no fluxo de alteração de categoria, o campo indica a categoria 
de destino; nos fluxos de recurso e de atualização de pontuação, permanece 
nulo. A abordagem foi adotada em vez da herança de tabela por subtipo 
(\textit{joined table inheritance}) por não haver, no escopo do módulo, 
atributos exclusivos de um tipo que justificassem tabelas próprias.

O \textbf{estado de rascunho} é consequência direta dessa modelagem. O 
estado \texttt{CREATED} representa uma solicitação em elaboração, que pode 
ser salva e editada antes da submissão. Nesse estado, o campo 
\texttt{id\_category} pode ser nulo, permitindo que o usuário salve a 
solicitação parcialmente preenchida e escolha o tipo mais adiante. A 
submissão para análise corresponde à transição de \texttt{CREATED} para 
\texttt{PENDING}, momento a partir do qual o tipo passa a ser obrigatório 
na prática, ainda que a restrição não esteja expressa no esquema. A 
alternativa — tornar \texttt{id\_category} obrigatório e criar uma tabela 
separada para rascunhos — foi descartada por duplicar a estrutura e a 
lógica de consulta sem ganho proporcional.

A \textbf{dupla representação de estados e categorias} — enumeração em 
Python e tabela de referência no banco — atende a duas necessidades 
distintas. A enumeração serve à aplicação, garantindo, em tempo de 
desenvolvimento, que apenas valores conhecidos sejam usados e permitindo 
comparações explícitas no código. A tabela de referência serve ao banco, 
garantindo integridade referencial e viabilizando junções e filtros 
eficientes nas consultas de listagem. O custo é a disciplina de manter os 
dois em sincronia a cada novo valor; o benefício é que o estado de uma 
solicitação nunca fica armazenado como texto livre, o que facilitaria 
inconsistências e inviabilizaria restrições declarativas.

A \textbf{chave estrangeira para \texttt{users}}, e não para a entidade de 
especialista, reflete no esquema o desafio identificado na 
Seção~\ref{sec:regras-negocio}: a categoria do especialista está em 
\texttt{users}, e não em \texttt{specialists}. A solicitação referencia 
o usuário e, quando aplicável, indica a categoria de destino em 
\texttt{specialist\_type\_id}, sem duplicar a classificação corrente. Isso 
mantém o modelo coerente com o restante do sistema e evita que a 
solicitação armazene estados intermediários que precisariam ser 
sincronizados com a tabela de origem.

A \textbf{indexação} foi definida a partir das consultas efetivamente 
implementadas nas rotas (Seção~\ref{sec:endpoints}): listagem por usuário 
(\texttt{idx\_solicitation\_user}), listagem por situação 
(\texttt{idx\_solicitation\_status}) e ordenação por data de criação 
(\texttt{idx\_solicitation\_created}). Consultas por identificador já são 
otimizadas pela chave primária e não exigem índice adicional. A escolha 
privilegia os caminhos de leitura da interface — listagem do solicitante, 
listagem do avaliador e ordenação cronológica — em vez de uma indexação 
genérica e não fundamentada.

O \textbf{armazenamento do arquivo} foi delegado ao sistema de arquivos: 
o atributo \texttt{file\_path} guarda o caminho do documento anexado, não 
o conteúdo binário. Essa decisão mantém o banco fora do caminho de 
operações de I/O de arquivos, que são executadas fora da transação 
principal, e evita o inchaço da tabela. A referência ao arquivo permanece 
no mesmo registro da solicitação, mantendo a consistência entre o 
documento e os metadados que o descrevem.

Por fim, o \textbf{controle temporal} foi delegado ao banco: 
\texttt{created\_at} recebe o instante corrente na inserção e 
\texttt{last\_update} é atualizado automaticamente a cada alteração, de 
modo que o momento da última modificação é registrado mesmo em operações 
que não passem pela aplicação. Isso atende parcialmente ao RNF02, no 
aspecto temporal da rastreabilidade; a dimensão semântica — quem alterou, 
de qual estado para qual estado, com qual justificativa — é atendida pelo 
mecanismo de snapshot e pela linha do tempo, descritos na 
Seção~\ref{sec:snapshot-auditoria}. A atomicidade das operações (RNF01), 
por sua vez, não é garantida no modelo, mas na camada de serviço, por 
meio da unidade de trabalho do SQLAlchemy (\texttt{session}): o modelo 
define as restrições; a transação define o momento em que elas são 
verificadas em conjunto.